\section{三角函数的图像}
\subsection{正弦函数、余弦函数的图像与性质}
\clearpage
\subsection{\texorpdfstring{$A\sin(\omega x+\varphi)+h$}s的图像与性质}%\texorpdfstring命令后面要加一个被吞掉的字符
在上一节中详细地介绍了\(\sin x \)与\(\cos x\)的图像，接下来就可以将过去学习的伸缩变换、平移变换作用到三角函数上，看看会有什么新的变化。
\begin{example}
    在同一直角坐标系中绘制\(y=\sin ,y=2\sin x\)的图像。
\end{example}
\begin{example}
    在同一直角坐标系中绘制\(y=\sin x,y=\sin 2x\)的图像。
\end{example}
\begin{example}
    在同一直角坐标系中绘制\(y=\sin x,y=\sin (x+\dfrac{\pi}{3})\)的图像。
\end{example}
\begin{example}
    在同一直角坐标系中绘制\(y=\sin x,y=\sin x+1\)的图像。
\end{example}
\begin{example}
    在同一直角坐标系中绘制\(y=\sin x,y=\sin(2x+\dfrac{\pi}{3})\)的图像。
\end{example}

\begin{example}
    \(f(x)=\sin(\omega x +\dfrac{\pi}{4})\)\((\omega>0)\)在\((\dfrac{\pi}{2},\dfrac{3\pi}{2})\)内没有最值，
    求\(\omega\)的取值范围。
\end{example}


\begin{example}[2022年全国甲卷]
    设函数 \( f(x) = \sin(\omega x + \dfrac{\pi}{3}) \) 在区间 \( (0, \pi) \) 恰有三个极值点、两个零点，求 \(\omega\) 的取值范围
\end{example}

\begin{example}[2016年新课标I卷]
    已知函数 \( f(x) = \sin(\omega x + \varphi) \left(\omega > 0, \ |\varphi| \leq \dfrac{\pi}{2}\right) \)，
    \( x = -\dfrac{\pi}{4} \) 为 \( f(x) \) 的零点，\( x = \dfrac{\pi}{4} \) 为 
    \( y = f(x) \) 图象的对称轴，且 \( f(x) \) 在 \( (\dfrac{\pi}{18}, \dfrac{5\pi}{36}) \) 上单调，求\(\omega\) 的最大值。
\end{example}


\begin{example}
已知函数 \( f(x) = 2\sin(\omega x + \varphi) \) (\(\omega > 0\)) 的图象关于直线 \( x = \dfrac{\pi}{3} \) 对称, 且 \( f\left(\dfrac{\pi}{12}\right) = 0 \), 求 \(\omega\) 的最小值为

\end{example}
\begin{exercise}
已知函数 \( f(x) = A\sin(\omega x + \varphi) \) (\(A > 0, \omega > 0, |\varphi| \leq \dfrac{\pi}{2}\))，满足 \( f\left(-\dfrac{\pi}{6}\right) = 0 \) 且 \(\forall x \in \mathbb{R} \) 都有 \( f(x) = f\left(\dfrac{2\pi}{3} - x\right) \)，若 \( f(x) \) 在 \( \left(\dfrac{5\pi}{36}, \dfrac{2\pi}{9}\right) \) 上单调，求\(\omega\) 的最大值为
\end{exercise}






















\clearpage
\subsection{正切函数的图像与性质}
\clearpage
\subsection{\texorpdfstring{$A\tan(\omega x+\varphi)+h$}s的图像与性质}
\clearpage
\subsection{余切、正割、余割函数的图像与性质**}
\clearpage
\subsection{反三角函数的图像与性质**}
\clearpage
\review





\clearpage
\hw
\begin{homework}[2012年新课标卷]
    已知 \(\omega > 0\)，函数 \( f(x) = \sin(\omega x + \dfrac{\pi}{4}) \) 在区间 \((\dfrac{\pi}{2}, \pi)\) 上单调递减，则实数 \(\omega\) 的取值范围是（\quad  ）

\begin{table}[H]
    \centering
\begin{tabular*}{\hsize}{@{}@{\extracolsep{\fill}}llll@{}}
   \(\text{A.}  \left[\dfrac{1}{2}, \dfrac{5}{4}\right] \)  & \( \text{B.}  \left[\dfrac{1}{2}, \dfrac{3}{4}\right] \) & \(    \text{C.} \left(0, \dfrac{1}{2}\right] \) & $\text{D.}  (0, 2]$
\end{tabular*}
\end{table}

\end{homework}

\begin{homework}
函数 \( f(x) = \sin\left(\omega x + \dfrac{\pi}{3}\right) \)   $ (\omega > 0)$ 在 \( \left(0, \dfrac{\pi}{3}\right) \) 单调递增，\(\omega\) 的取值范围为\_\_\_\_\_。
\end{homework}

\begin{homework}[2016天津卷]
    已知函数 \( f(x) = \sin\left(\omega x - \dfrac{\pi}{4}\right) (\omega > 0) \)，若 \( f(x) \) 在区间 \( (\pi, 2\pi) \) 内没有零点，则 \(\omega\) 的取值范围是（\quad ）

\begin{table}[H]
    \centering
\begin{tabular*}{\hsize}{@{}@{\extracolsep{\fill}}llll@{}}
  $\text{A.}  \left(0, \dfrac{1}{8}\right] $ & $\text{B.}\left(0, \dfrac{1}{4}\right] \cup \left[\dfrac{5}{8}, 1\right) $ & $\text{C.} \left(0, \dfrac{5}{8}\right] $&$\text{D.} (0, 8] \cup \left[\dfrac{1}{4}, \dfrac{5}{8}\right]$
\end{tabular*}
\end{table}

\end{homework}
\begin{homework}
  函数 \( f(x) = \sin\left(\omega x - \frac{\pi}{6}\right) \) $( \omega > 0)$ 在 \( (0, 2\pi) \) 有且仅有一个零点，则 \(\omega\) 的取值范围为\_\_\_\_\_。
\end{homework}
\begin{homework}
函数 \( f(x) = \cos\left(\omega x + \frac{\pi}{6}\right) \)   $( \omega > 0)$ 在 \( \left(-\frac{\pi}{3}, \frac{\pi}{2}\right) \) 有且仅有$ 1$ 个零点，\(\omega\) 的取值范围为\_\_\_\_\_。
\end{homework}
\begin{homework}
    函数 \( f(x) = \sin\left(\omega x + \frac{\pi}{6}\right) \) 在 \( \left(-\frac{\pi}{2}, 0\right) \) 没有零点，\(\omega\) 的取值范围为\_\_\_\_\_。
\end{homework}
\begin{homework}[2018北京卷理]
 设函数 \( f(x) = \cos\left(\omega x - \dfrac{\pi}{6}\right) \) (\(\omega > 0\))。若 \( f(x) \leq f\left(\frac{\pi}{4}\right) \) 对任意的实数 \( x \) 都成立，则 \(\omega\) 的最小值为\_\_\_\_\_。
\end{homework}
\begin{homework}
    已知函数 \( f(x) = 6\cos(\omega x + \varphi) (\omega > 0, 0 < \varphi < \pi) \)，\( \forall x \in \mathbb{R} \) 都有 \( f(x) \leq |f\left(\dfrac{\pi}{6}\right)| \)，且 \( x = -\dfrac{\pi}{6} \) 是 \( f(x) \) 的一个零点。若 \( g(x) = f(x) - 6 \) 在 \( \left(\dfrac{\pi}{15}, \dfrac{\pi}{6}\right) \) 上有且只有一个零点，则 \(\omega\) 的最大值为\_\_\_\_\_。
\end{homework}
\begin{homework}
    
\end{homework}
\begin{homework}
    
\end{homework}
\begin{homework}
    
\end{homework}
\begin{homework}
    
\end{homework}
\begin{homework}
    
\end{homework}
\begin{homework}
    
\end{homework}
\begin{homework}
    
\end{homework}
\begin{homework}
    
\end{homework}
\begin{homework}
    
\end{homework}
\begin{homework}
    
\end{homework}




\clearpage
